\chapter{Chromatin Accessibility Profiling}
\label{ch:chromatin}

\begin{keyconcept}
Chromatin accessibility refers to the degree to which nuclear DNA is physically accessible to regulatory proteins such as transcription factors, RNA polymerases, and chromatin remodelers. In eukaryotic cells, DNA is packaged around histone octamers into nucleosomes, and the local density and positioning of these nucleosomes determines whether a genomic region is ``open'' (accessible) or ``closed'' (compacted). Open chromatin regions typically correspond to active regulatory elements---promoters, enhancers, silencers, and insulators. Sequencing-based methods such as ATAC-seq and DNase-seq allow researchers to map these accessible regions genome-wide, providing a snapshot of the regulatory landscape of any cell type.
\end{keyconcept}

%% ============================================================
\section{Why Chromatin Accessibility Matters}
\label{sec:chromatin:why}

The human genome contains approximately 3.2 billion base pairs, yet only a small fraction of this sequence is accessible in any given cell type. The accessible fraction---typically 2--4\% of the genome---encompasses the active \textbf{cis-regulatory elements} that control gene expression:

\begin{itemize}[itemsep=3pt]
  \item \textbf{Active promoters:} Regions immediately upstream of transcription start sites (TSSs) where the transcription machinery assembles.
  \item \textbf{Enhancers:} Distal regulatory elements that boost transcription of target genes, often located tens to hundreds of kilobases away.
  \item \textbf{Insulators/boundary elements:} Regions that prevent inappropriate enhancer--promoter interactions, often bound by CTCF.
  \item \textbf{Silencers:} Elements that repress transcription of target genes.
  \item \textbf{Transcription factor (TF) binding sites:} Accessible DNA motifs where sequence-specific TFs bind to regulate gene expression.
\end{itemize}

Chromatin accessibility profiling provides a direct readout of regulatory element activity. Unlike histone modification mapping, which requires specific antibodies for each mark, accessibility assays capture \textit{all} open regulatory elements in a single experiment. Changes in accessibility between conditions (e.g., disease vs.\ healthy, differentiated vs.\ stem cells) reveal the regulatory programs that drive cell-state transitions.

%% ============================================================
\section{DNase-seq}
\label{sec:chromatin:dnase}

\subsection{Principle}
\label{subsec:chromatin:dnase-principle}

DNase I is an endonuclease that preferentially cleaves DNA that is not protected by nucleosomes or tightly bound proteins. Regions of open chromatin are hypersensitive to DNase I digestion, and these \textbf{DNase I hypersensitive sites (DHSs)} have been used since the 1980s to identify regulatory elements.

In \textbf{DNase-seq}:
\begin{enumerate}[itemsep=2pt]
  \item Intact nuclei are isolated from cells.
  \item Nuclei are treated with a controlled amount of DNase I enzyme.
  \item The small DNA fragments released from hypersensitive sites are size-selected (typically $<$ 300\basepair{}).
  \item Fragments are adapter-ligated and sequenced.
  \item Reads pile up at DHSs, producing peaks in the coverage signal.
\end{enumerate}

\subsection{The ENCODE Project and DHS Mapping}
\label{subsec:chromatin:encode}

DNase-seq was the method of choice for the \acrfull{ENCODE} project, which mapped DHSs in hundreds of human cell types and tissues. The ENCODE DHS maps identified approximately 2.9 million DHSs across the human genome, revealing that $\sim$40\% of the genome is accessible in at least one cell type, even though only 2--3\% is accessible in any individual cell type. This work demonstrated the vast regulatory diversity of human cells and provided a foundational resource for interpreting non-coding genetic variants.

\begin{tipbox}
While DNase-seq has been largely superseded by ATAC-seq for new experiments (due to ATAC-seq's simpler protocol and lower cell requirements), the massive DNase-seq datasets generated by ENCODE remain invaluable reference resources. The ENCODE portal (\url{https://www.encodeproject.org}) provides freely accessible DHS maps for hundreds of cell types.
\end{tipbox}

%% ============================================================
\section{ATAC-seq}
\label{sec:chromatin:atac}

\subsection{Principle: Tn5 Transposase-Mediated Tagmentation}
\label{subsec:chromatin:atac-principle}

\acrfull{ATAC-seq}, introduced by Buenrostro et al.\ in 2013, has become the dominant method for chromatin accessibility profiling. The key innovation is the use of a hyperactive \textbf{Tn5 transposase} loaded with sequencing adapters. The Tn5 transposase simultaneously fragments DNA and inserts sequencing adapters into accessible chromatin regions---a process called \textbf{tagmentation}:

\begin{enumerate}[itemsep=2pt]
  \item The Tn5 transposase can only access DNA that is not wrapped around nucleosomes or occluded by other proteins.
  \item When it inserts into accessible DNA, it creates a 9\basepair{} staggered cut and ligates sequencing adapters to both fragment ends.
  \item Nucleosome-protected regions are refractory to Tn5 insertion.
  \item After tagmentation, DNA is purified and amplified by PCR with primers complementary to the inserted adapters.
  \item Sequencing reveals the locations where Tn5 inserted, which correspond to accessible chromatin.
\end{enumerate}

\begin{figure}[htbp]
\centering
\begin{tikzpicture}[
  scale=0.95,
  nuc/.style={circle, draw=MidnightBlue!70, fill=MidnightBlue!15, minimum size=0.8cm, font=\tiny\bfseries},
  dna/.style={very thick, BrickRed!70},
  tn5/.style={isosceles triangle, draw=OliveGreen!80, fill=OliveGreen!20, minimum size=0.4cm, shape border rotate=90, inner sep=1pt, font=\tiny\bfseries},
  adapter/.style={very thick, Goldenrod},
  arrow/.style={-{Stealth[length=2.5mm]}, thick, MidnightBlue},
  label/.style={font=\scriptsize}
]

% Title
\node[font=\bfseries\small, MidnightBlue] at (6, 7.8) {ATAC-seq: Tn5 Transposition into Open Chromatin};

% --- Panel A: Chromatin with nucleosomes ---
\node[font=\scriptsize\bfseries, MidnightBlue] at (-1.5, 6.8) {A};

% Draw DNA as a line through nucleosomes
\draw[dna] (0,6.5) -- (1.5,6.5);
% Nucleosome 1
\draw[dna] (1.5,6.5) arc (-90:270:0.45);
\node[nuc] at (1.5,6.95) {Nuc};
\draw[dna] (1.95,6.5) -- (3.2,6.5);
% Open region
\draw[dna, OliveGreen!70] (3.2,6.5) -- (5.8,6.5);
\node[font=\scriptsize, OliveGreen!80, above] at (4.5,6.6) {Open chromatin};
% Nucleosome 2
\draw[dna] (5.8,6.5) -- (6.5,6.5);
\draw[dna] (6.5,6.5) arc (-90:270:0.45);
\node[nuc] at (6.5,6.95) {Nuc};
\draw[dna] (6.95,6.5) -- (8.2,6.5);
% Open region 2
\draw[dna, OliveGreen!70] (8.2,6.5) -- (9.8,6.5);
% Nucleosome 3
\draw[dna] (9.8,6.5) -- (10.5,6.5);
\draw[dna] (10.5,6.5) arc (-90:270:0.45);
\node[nuc] at (10.5,6.95) {Nuc};
\draw[dna] (10.95,6.5) -- (12,6.5);

% --- Panel B: Tn5 insertion ---
\node[font=\scriptsize\bfseries, MidnightBlue] at (-1.5, 4.8) {B};

\draw[dna] (0,4.5) -- (1.5,4.5);
\draw[dna] (1.5,4.5) arc (-90:270:0.45);
\node[nuc] at (1.5,4.95) {Nuc};
\draw[dna] (1.95,4.5) -- (3.0,4.5);

% Tn5 insertions in open region
\node[tn5] (tn5a) at (3.5,5.2) {\textcolor{OliveGreen!80}{Tn5}};
\draw[OliveGreen!60, thick, -{Stealth}] (tn5a.south) -- (3.5,4.7);
\node[tn5] (tn5b) at (5.5,5.2) {\textcolor{OliveGreen!80}{Tn5}};
\draw[OliveGreen!60, thick, -{Stealth}] (tn5b.south) -- (5.5,4.7);

\draw[dna, OliveGreen!70] (3.0,4.5) -- (6.0,4.5);

\draw[dna] (6.0,4.5) -- (6.5,4.5);
\draw[dna] (6.5,4.5) arc (-90:270:0.45);
\node[nuc] at (6.5,4.95) {Nuc};
\draw[dna] (6.95,4.5) -- (8.2,4.5);

% Blocked Tn5 at nucleosome
\node[tn5, fill=gray!20, draw=gray] (tn5c) at (1.5,5.8) {\textcolor{gray}{Tn5}};
\draw[gray, thick] (1.2,5.5) -- (1.8,5.5);
\node[font=\tiny, gray, above] at (1.5,5.9) {blocked};

% --- Panel C: Resulting fragments ---
\node[font=\scriptsize\bfseries, MidnightBlue] at (-1.5, 2.8) {C};

% NFR fragment (short)
\draw[adapter] (1,2.8) -- (1.8,2.8);
\draw[dna] (1.8,2.8) -- (3.8,2.8);
\draw[adapter] (3.8,2.8) -- (4.6,2.8);
\node[font=\tiny, below] at (2.8,2.6) {NFR fragment ($<$ 150 bp)};

% Mono-nucleosome fragment
\draw[adapter] (6,2.8) -- (6.8,2.8);
\draw[dna] (6.8,2.8) -- (7.3,2.8);
\draw[dna] (7.3,2.8) arc (-90:270:0.35);
\node[nuc, minimum size=0.55cm] at (7.3,3.15) {\tiny Nuc};
\draw[dna] (7.65,2.8) -- (8.5,2.8);
\draw[adapter] (8.5,2.8) -- (9.3,2.8);
\node[font=\tiny, below] at (7.9,2.6) {Mono-nuc fragment ($\sim$ 200 bp)};

% Legend
\draw[adapter, thick] (0.5,1.5) -- (1.5,1.5);
\node[font=\tiny, right] at (1.6,1.5) {Sequencing adapter (inserted by Tn5)};
\draw[dna, thick] (5.5,1.5) -- (6.5,1.5);
\node[font=\tiny, right] at (6.6,1.5) {Genomic DNA};

% Fragment size annotation
\draw[arrow] (0, 0.5) -- (12, 0.5);
\node[font=\scriptsize, above] at (6, 0.6) {Fragment Size Distribution};
\node[font=\tiny, below] at (2.5, 0.3) {NFR ($<$150 bp)};
\node[font=\tiny, below] at (5.5, 0.3) {Mono-nuc ($\sim$200 bp)};
\node[font=\tiny, below] at (8.5, 0.3) {Di-nuc ($\sim$400 bp)};

\end{tikzpicture}
\caption[ATAC-seq principle: Tn5 transposase tagmentation of open chromatin.]{ATAC-seq principle. (A) Chromatin consists of nucleosome-wrapped regions (closed) and nucleosome-free regions (open). (B) The Tn5 transposase, pre-loaded with sequencing adapters, preferentially inserts into accessible (open) regions. Nucleosome-bound DNA blocks Tn5 insertion. (C) The resulting fragments include short nucleosome-free region (NFR) fragments and longer mono-nucleosome fragments, producing a characteristic fragment size distribution with peaks at $<$ 150\basepair{} (NFR), $\sim$200\basepair{} (mono-nucleosomal), and $\sim$400\basepair{} (di-nucleosomal).}
\label{fig:chromatin:atac-principle}
\end{figure}

\subsection{Standard ATAC-seq Protocol}
\label{subsec:chromatin:atac-protocol}

\begin{protocolbox}[title={Standard ATAC-seq Protocol (Omni-ATAC)}]
\begin{enumerate}[itemsep=2pt]
  \item \textbf{Cell preparation:} Harvest 50,000--100,000 cells. Count cells and assess viability ($>$90\% recommended).
  \item \textbf{Lysis:} Gently lyse cells with a non-ionic detergent (e.g., 0.1\% NP-40, 0.1\% Tween-20, 0.01\% digitonin in Omni-ATAC) to release nuclei while keeping nuclei intact.
  \item \textbf{Tagmentation:} Incubate nuclei with Tn5 transposase (e.g., Illumina Tagment DNA Enzyme) at 37\textdegree{}C for 30 minutes. The enzyme simultaneously fragments and adapter-tags accessible DNA.
  \item \textbf{DNA purification:} Clean up tagmented DNA using a column or SPRI beads.
  \item \textbf{PCR amplification:} Amplify library using indexed primers (5--12 cycles, determined by qPCR to avoid over-amplification).
  \item \textbf{Size selection (optional):} Remove large fragments ($>$1000\basepair{}) by bead-based selection. Some protocols skip this.
  \item \textbf{QC:} Assess fragment size distribution on Bioanalyzer/TapeStation. A successful ATAC-seq library shows the characteristic nucleosomal laddering pattern.
  \item \textbf{Sequencing:} Paired-end 50 or 75\basepair{} reads. 50--75 million read pairs per sample recommended.
\end{enumerate}
\end{protocolbox}

\subsection{ATAC-seq Protocol Variants}
\label{subsec:chromatin:atac-variants}

Several optimised ATAC-seq protocols have been developed for specific applications:

\begin{itemize}[itemsep=3pt]
  \item \textbf{Omni-ATAC:} The recommended standard protocol. Uses a combination of detergents (NP-40 + Tween-20 + digitonin) for more complete lysis, resulting in cleaner signal and reduced mitochondrial contamination. Works across diverse cell types.
  \item \textbf{Fast-ATAC:} Optimised for blood cells and other suspension cells. Cells are tagmented directly without a separate lysis step (Tn5 + detergent in a single reaction). Reduces hands-on time and improves results for difficult cell types.
  \item \textbf{miniATAC-seq:} Designed for very low cell numbers (500--5,000 cells). Uses modifications to reduce material loss at each step.
  \item \textbf{ATAC-seq on tissues:} Fresh or flash-frozen tissues are dissociated into single nuclei (mechanical dissociation or enzymatic digestion), then standard ATAC-seq is performed. snATAC-seq on nuclei avoids the need for live cells.
  \item \textbf{mtscATAC-seq:} Modified protocol that retains mitochondrial DNA information alongside chromatin accessibility in single cells, enabling clonal tracking based on mitochondrial mutations.
\end{itemize}

\subsection{Tn5 Insertion Bias}
\label{subsec:chromatin:tn5-bias}

The Tn5 transposase has a modest sequence preference for insertion. This \textbf{Tn5 bias} can affect footprinting analysis (where exact cut-site positions are used to infer TF binding). Bias correction methods include:

\begin{itemize}[itemsep=2pt]
  \item \textbf{Naked DNA control:} Tagment purified genomic DNA (no chromatin) with Tn5 to empirically measure insertion bias.
  \item \textbf{Computational correction:} Tools like TOBIAS and seqOutBias model and correct for the Tn5 sequence preference.
  \item \textbf{k-mer-based normalisation:} Normalise cut-site frequencies by the local sequence composition.
\end{itemize}

\begin{warningbox}
A major source of wasted sequencing in ATAC-seq is \textbf{mitochondrial reads}. Because mitochondrial DNA is abundant and lacks histones (and is therefore fully accessible), mitochondrial reads can constitute 30--60\% of a standard ATAC-seq library. The Omni-ATAC protocol reduces this to $\sim$5--15\%, but filtering mitochondrial reads remains an essential bioinformatics step. Always report the fraction of mitochondrial reads as a QC metric.
\end{warningbox}

%% ============================================================
\section{Single-Cell ATAC-seq (scATAC-seq)}
\label{sec:chromatin:scatac}

\subsection{Technology and Platforms}
\label{subsec:chromatin:scatac-platforms}

Single-cell ATAC-seq (scATAC-seq) profiles chromatin accessibility in individual cells, revealing cell-type-specific regulatory landscapes within complex tissues. The dominant commercial platform is the \textbf{10x Genomics Chromium Single Cell ATAC} system:

\begin{enumerate}[itemsep=2pt]
  \item Nuclei are isolated and tagmented in bulk with Tn5.
  \item Tagmented nuclei are partitioned into nanoliter-scale gel bead-in-emulsion (GEM) droplets, each containing a single nucleus and a barcoded gel bead.
  \item Within each GEM, linear amplification attaches cell-specific barcodes to the tagmented fragments.
  \item After breaking the emulsions, all barcoded fragments are pooled, amplified, and sequenced.
  \item Demultiplexing by barcode assigns fragments to individual cells.
\end{enumerate}

Other platforms include \textbf{sci-ATAC-seq} (combinatorial indexing without droplets, enabling very high cell throughput), \textbf{SHARE-seq} (simultaneous ATAC + RNA from the same cell), and \textbf{10x Multiome} (combined scATAC + scRNA-seq from the same nucleus).

\subsection{The Sparsity Challenge}
\label{subsec:chromatin:scatac-sparsity}

scATAC-seq data is inherently \textbf{extremely sparse}. Each cell has only two copies of each autosomal locus (diploid genome), and the detection efficiency is far below 100\%. A typical scATAC-seq experiment recovers only 5,000--20,000 unique fragments per cell. The resulting cell-by-peak matrix is $>$95\% zeros, making analysis fundamentally different from scRNA-seq:

\begin{itemize}[itemsep=2pt]
  \item Individual peaks in single cells are essentially \textbf{binary} (0 or 1 fragment).
  \item Clustering and dimensionality reduction must use methods suited to binary/sparse data.
  \item Feature aggregation (peaks $\rightarrow$ gene activity scores, or peaks $\rightarrow$ TF motif scores) is essential to accumulate signal.
\end{itemize}

\subsection{Analysis Tools and Frameworks}
\label{subsec:chromatin:scatac-tools}

\begin{itemize}[itemsep=3pt]
  \item \textbf{ArchR:} Comprehensive R package for scATAC-seq analysis. Uses Arrow files for memory-efficient processing. Implements iterative LSI (latent semantic indexing) for dimensionality reduction, peak calling with MACS2, gene activity scores, motif enrichment, and trajectory analysis.
  \item \textbf{Signac:} Extension of Seurat for scATAC-seq. Provides integration with scRNA-seq data, chromVAR motif analysis, gene activity scoring, and multimodal analysis.
  \item \textbf{SnapATAC2:} Python-based tool using an anndata framework. Fast and memory-efficient; supports spectral embedding, cell clustering, and peak calling.
  \item \textbf{cisTopic:} Uses Latent Dirichlet Allocation (topic modelling) to identify co-accessible regulatory modules. Each ``topic'' represents a set of genomic regions that tend to be accessible together, often corresponding to specific cell types or regulatory programs.
\end{itemize}

Key analysis steps in scATAC-seq:
\begin{enumerate}[itemsep=2pt]
  \item \textbf{Cell filtering:} Remove low-quality cells based on unique fragment count, TSS enrichment, and nucleosome signal.
  \item \textbf{Dimensionality reduction:} LSI (TF-IDF normalisation $\rightarrow$ SVD) or topic modelling.
  \item \textbf{Clustering:} Graph-based clustering (Louvain, Leiden) on the LSI embedding.
  \item \textbf{Peak calling:} Call peaks per cluster (pseudo-bulk) using MACS2/MACS3.
  \item \textbf{Gene activity scores:} Estimate gene activity by summing accessibility signal in gene body and promoter regions.
  \item \textbf{Motif enrichment:} Identify enriched TF motifs in cluster-specific peaks using chromVAR or HOMER.
  \item \textbf{Integration with scRNA-seq:} Transfer cell-type labels from scRNA-seq to scATAC-seq using integration methods (e.g., Seurat WNN, ArchR gene activity integration).
\end{enumerate}

\begin{tipbox}
When designing a scATAC-seq experiment, aim for at least 25,000 unique fragments per cell. Libraries with $<$5,000 fragments per cell are typically too sparse for meaningful analysis. Use TSS enrichment ($>$4 is good; $>$8 is excellent) as the primary QC metric for library quality.
\end{tipbox}

%% ============================================================
\section{Other Accessibility and Nucleosome Methods}
\label{sec:chromatin:other}

\subsection{FAIRE-seq}
\label{subsec:chromatin:faire}

Formaldehyde-Assisted Isolation of Regulatory Elements (FAIRE-seq) was an early accessibility method. Cells are crosslinked with formaldehyde, then DNA is sheared and subjected to phenol-chloroform extraction. Protein-free (accessible) DNA partitions into the aqueous phase, while protein-bound (nucleosomal) DNA partitions into the organic phase. FAIRE-seq is now rarely used due to high background and replacement by ATAC-seq.

\subsection{MNase-seq}
\label{subsec:chromatin:mnase}

Micrococcal nuclease (MNase) digests linker DNA between nucleosomes but is impeded by histone-bound DNA. By sequencing the protected fragments:

\begin{itemize}[itemsep=2pt]
  \item \textbf{Nucleosome positioning maps} are generated, showing where nucleosomes sit along the genome.
  \item Well-positioned nucleosomes flanking TSSs and TF binding sites are characteristic features of active genes.
  \item MNase-seq complements ATAC-seq: ATAC-seq maps open regions, while MNase-seq maps nucleosome positions.
\end{itemize}

\subsection{NOMe-seq}
\label{subsec:chromatin:nome}

Nucleosome Occupancy and Methylome Sequencing (NOMe-seq) provides \textbf{simultaneous measurement of chromatin accessibility and DNA methylation on the same DNA molecule}:

\begin{enumerate}[itemsep=2pt]
  \item Treat intact nuclei with GpC methyltransferase (M.CviPI), which methylates cytosines in GpC dinucleotides that are accessible.
  \item Nucleosome-protected GpC sites are not methylated by M.CviPI.
  \item Bisulfite sequencing then reads both: (a) endogenous CpG methylation (from the cell's own methylation), and (b) exogenous GpC methylation (marking accessible chromatin).
  \item On long reads (nanopore or PacBio), this provides single-molecule, phased measurements of both methylation and accessibility.
\end{enumerate}

%% ============================================================
\section{Comprehensive Comparison of Accessibility Methods}
\label{sec:chromatin:comparison}

\begin{table}[htbp]
\centering
\caption[Comparison of chromatin accessibility profiling methods.]{Comparison of major chromatin accessibility profiling approaches.}
\label{tab:chromatin:comparison}
\small
\begin{tabularx}{\textwidth}{l c c c c X}
\toprule
\textbf{Method} & \textbf{Input Cells} & \textbf{Reads Needed} & \textbf{Resolution} & \textbf{Time} & \textbf{Key Features} \\
\midrule
DNase-seq & $10^{6}$--$10^{7}$ & 30--50M & $\sim$1\basepair{} (cut site) & 2--3 days & ENCODE standard; historical \\
\addlinespace
ATAC-seq & 50K--100K & 50--75M & $\sim$1\basepair{} (cut site) & $\sim$3 hours & Simple; dominant method \\
\addlinespace
Omni-ATAC & 50K--100K & 50--75M & $\sim$1\basepair{} & $\sim$3 hours & Improved ATAC; less mito contamination \\
\addlinespace
scATAC-seq & 5K--10K nuclei & 20--50K/cell & Single cell & $\sim$4--5 hours & Cell-type-resolved accessibility \\
\addlinespace
FAIRE-seq & $10^{6}$--$10^{7}$ & 20--30M & $\sim$200\basepair{} & 1--2 days & Superseded; no enzyme needed \\
\addlinespace
MNase-seq & $10^{6}$--$10^{7}$ & 100--200M & $\sim$10\basepair{} & 1--2 days & Nucleosome positions \\
\addlinespace
NOMe-seq & $10^{5}$--$10^{6}$ & 100--200M (WG) & Single molecule & 1--2 days & Accessibility + methylation jointly \\
\bottomrule
\end{tabularx}
\end{table}

%% ============================================================
\section{Bioinformatics Workflow for ATAC-seq}
\label{sec:chromatin:bioinformatics}

\subsection{Analysis Overview}
\label{subsec:chromatin:analysis-overview}

The standard ATAC-seq bioinformatics pipeline includes the following steps:

\begin{enumerate}[itemsep=2pt]
  \item \textbf{Quality control:} FastQC on raw reads. Check for adapter contamination, insert size distribution.
  \item \textbf{Adapter trimming:} Trim Galore or fastp to remove Nextera adapter sequences.
  \item \textbf{Alignment:} Bowtie2 in paired-end mode with specific settings (e.g., \texttt{-X 2000} to allow large insert sizes, \texttt{--very-sensitive}).
  \item \textbf{Filtering:}
    \begin{itemize}[itemsep=1pt]
      \item Remove mitochondrial reads (\texttt{samtools view} excluding chrM).
      \item Remove PCR duplicates (Picard MarkDuplicates or samtools markdup).
      \item Remove low-quality alignments (MAPQ $<$ 30).
      \item Remove reads in ENCODE blacklist regions.
      \item For paired-end data, keep only properly paired reads.
    \end{itemize}
  \item \textbf{Tn5 offset adjustment:} Shift reads by +4/$-$5\basepair{} to centre on the Tn5 cut site (accounts for the 9\basepair{} duplication).
  \item \textbf{Peak calling:} MACS2 or MACS3 with appropriate settings (\texttt{--nomodel --shift -100 --extsize 200} for Tn5 cut-site-centred calling, or using the paired-end BAMPE mode).
  \item \textbf{Quality metrics:}
    \begin{itemize}[itemsep=1pt]
      \item Fragment size distribution (nucleosomal laddering).
      \item TSS enrichment score.
      \item FRiP (Fraction of Reads in Peaks).
    \end{itemize}
  \item \textbf{Differential accessibility:} DiffBind or DESeq2 on a consensus peak set.
  \item \textbf{Motif analysis:} HOMER \texttt{findMotifsGenome.pl}, chromVAR, MEME Suite.
  \item \textbf{Footprinting:} TOBIAS or HINT-ATAC for TF footprint detection.
\end{enumerate}

\subsubsection{MACS2 Dynamic Poisson Model for Peak Calling}
\label{subsubsec:chromatin:macs2-poisson}

\software{MACS2} (Model-based Analysis of ChIP-Seq), the most widely used peak caller for both ChIP-seq and ATAC-seq data, employs a dynamic Poisson model that adapts to local background variation across the genome. Understanding this model clarifies how peaks are identified and why certain parameter choices matter.

\textbf{Global background estimation.} MACS2 first estimates a genome-wide average background rate $\lambda_{\text{genome}}$ by dividing the total number of reads by the effective genome size:
\begin{equation}
\label{eq:chromatin:macs2-global}
\lambda_{\text{genome}} = \frac{N_{\text{reads}}}{G_{\text{effective}}}
\end{equation}
where $N_{\text{reads}}$ is the total number of mapped reads and $G_{\text{effective}}$ is the mappable genome size (e.g., $2.7 \times 10^{9}$ for human).

\textbf{Local background estimation.} Because read density varies across the genome due to copy number variation, mappability, and genuine biological heterogeneity, a global rate alone is insufficient. MACS2 estimates local background rates using sliding windows of different sizes centred on each candidate peak:
\begin{itemize}[itemsep=2pt]
  \item $\lambda_{1k}$: background estimated from a 1\kilobase{} window around the peak.
  \item $\lambda_{5k}$: background estimated from a 5\kilobase{} window.
  \item $\lambda_{10k}$: background estimated from a 10\kilobase{} window.
\end{itemize}

The local lambda used for testing at each position is the maximum of all estimates, ensuring conservative peak calls:
\begin{equation}
\label{eq:chromatin:macs2-lambda}
\lambda_{\text{local}} = \max\!\left(\lambda_{1k},\; \lambda_{5k},\; \lambda_{10k},\; \lambda_{\text{genome}}\right)
\end{equation}

This dynamic approach prevents false-positive peak calls in regions with elevated background (e.g., copy number gains, accessible but non-regulatory regions).

\textbf{Significance testing.} Given the local background rate $\lambda_{\text{local}}$, MACS2 models the read count in a peak region as a Poisson random variable. The p-value for observing $k$ or more reads under the null (background) model is:
\begin{equation}
\label{eq:chromatin:macs2-pvalue}
p = P(X \geq k) = 1 - \sum_{i=0}^{k-1} \frac{\lambda_{\text{local}}^{i}\, e^{-\lambda_{\text{local}}}}{i!}
\end{equation}

\textbf{Multiple testing correction.} Because MACS2 tests millions of candidate positions across the genome, multiple testing correction is essential. MACS2 applies the \textbf{Benjamini--Hochberg (BH)} procedure to control the false discovery rate (FDR). The output q-value column reports BH-adjusted p-values, and the default threshold of $q < 0.05$ (or $q < 0.01$ for ATAC-seq) corresponds to controlling the FDR at 5\% (or 1\%).

\begin{tipbox}
For ATAC-seq data, MACS2 is typically run in \texttt{BAMPE} mode (paired-end), which uses the actual fragment boundaries rather than extending reads from a single end. This avoids the need for the \texttt{--nomodel --shift --extsize} parameters required for single-end data and provides more accurate peak boundaries. When using \texttt{BAMPE} mode, set \texttt{--keep-dup all} because PCR duplicates should already have been removed in the filtering step.
\end{tipbox}

\subsubsection{IDR: Irreproducible Discovery Rate}
\label{subsubsec:chromatin:idr}

The \textbf{Irreproducible Discovery Rate (IDR)} framework, developed for the ENCODE project, provides a principled statistical method for assessing the reproducibility of peak calls across biological replicates. IDR is particularly important for defining a set of high-confidence peaks.

\textbf{The copula mixture model.} IDR models the ranked peak signals from two replicates as arising from a mixture of two components:
\begin{enumerate}[itemsep=2pt]
  \item \textbf{Reproducible signal component:} True peaks that are consistently detected across replicates. The signals from the two replicates are correlated (modelled by a copula with high correspondence).
  \item \textbf{Irreproducible noise component:} Spurious peaks that appear in one replicate but not the other (or at very different ranks). The signals from the two replicates are essentially uncorrelated.
\end{enumerate}

Formally, for each peak $i$, let $(x_{i,1}, x_{i,2})$ be the signal intensities (e.g., $-\log_{10}(p)$) in replicates 1 and 2. The IDR model assumes:
\begin{equation}
\label{eq:chromatin:idr-mixture}
f(x_{i,1}, x_{i,2}) = \pi_1 \cdot f_{\text{repro}}(x_{i,1}, x_{i,2}) + (1 - \pi_1) \cdot f_{\text{irrepro}}(x_{i,1}, x_{i,2})
\end{equation}
where $\pi_1$ is the proportion of reproducible peaks, $f_{\text{repro}}$ is a bivariate distribution with high correlation (the reproducible component), and $f_{\text{irrepro}}$ is a bivariate distribution with low or zero correlation (the irreproducible component).

\textbf{IDR score.} For each peak, the IDR framework computes the posterior probability of belonging to the irreproducible component. The IDR score is the expected proportion of irreproducible discoveries at a given rank threshold. Peaks with \textbf{IDR $< 0.05$} are considered reproducible---analogous to an FDR threshold of 5\% but applied to replicate concordance rather than statistical significance.

\textbf{Practical use.} The ENCODE consortium recommends the following IDR workflow for ATAC-seq and ChIP-seq:
\begin{enumerate}[itemsep=2pt]
  \item Call peaks on each individual replicate using lenient thresholds (e.g., $p < 0.01$ without FDR correction).
  \item Run IDR on pairs of replicates to obtain a conservative, reproducible peak set (IDR $< 0.05$).
  \item Call peaks on the pooled data (all replicates merged) for the final peak set used in downstream analysis.
  \item Self-consistency analysis: split a single replicate into pseudo-replicates and verify that the IDR curve is well-behaved.
\end{enumerate}

\begin{warningbox}
IDR is designed for \textbf{sharp/narrow peaks} (such as those from TF ChIP-seq, H3K4me3, or ATAC-seq). It is not appropriate for broad marks (H3K27me3, H3K36me3) because broad peaks have fundamentally different signal characteristics that violate the IDR model assumptions. For broad marks, use overlap-based reproducibility metrics or the proportion of peaks shared between replicates.
\end{warningbox}

\subsubsection{Quality Metrics in Detail}

\textbf{Fragment Size Distribution:} A successful ATAC-seq library produces a characteristic ``nucleosomal ladder'' pattern when fragment sizes are plotted:
\begin{itemize}[itemsep=1pt]
  \item \textbf{Peak 1} ($<$ 150\basepair{}): nucleosome-free region (NFR) fragments---the primary signal of interest.
  \item \textbf{Peak 2} ($\sim$200\basepair{}): mono-nucleosomal fragments.
  \item \textbf{Peak 3} ($\sim$400\basepair{}): di-nucleosomal fragments.
  \item The ratio of NFR to mono-nucleosomal fragments indicates signal quality. Higher NFR enrichment is better.
\end{itemize}

\textbf{TSS Enrichment Score:} Calculated as the ratio of ATAC-seq signal at TSSs versus flanking regions ($\pm$ 2\kilobase{}). The ENCODE consortium recommends TSS enrichment $\geq$ 6 for acceptable quality and $\geq$ 10 for high quality.

\textbf{FRiP (Fraction of Reads in Peaks):} The proportion of reads that fall within called peaks. ENCODE recommends FRiP $\geq$ 0.3 (30\%). Low FRiP indicates high background or poor tagmentation.

\subsection{Snakemake Workflow}
\label{subsec:chromatin:snakemake}

\begin{lstlisting}[language=Python, caption={Snakemake workflow for ATAC-seq analysis.}, label={lst:chromatin:snakemake}]
# Snakefile for ATAC-seq analysis
configfile: "config.yaml"

SAMPLES = config["samples"]
GENOME = config["genome"]   # Bowtie2 index prefix
BLACKLIST = config["blacklist"]
EFFECTIVE_GENOME = config["effective_genome_size"]

rule all:
    input:
        "results/multiqc/multiqc_report.html",
        expand("results/peaks/{sample}_peaks.narrowPeak",
               sample=SAMPLES),
        expand("results/bigwig/{sample}.bw", sample=SAMPLES),
        "results/diffbind/differential_peaks.tsv",
        "results/motifs/homer_results/knownResults.html",
        expand("results/qc/{sample}_fragment_sizes.pdf",
               sample=SAMPLES)

rule trim_adapters:
    input:
        r1 = "raw_data/{sample}_R1.fastq.gz",
        r2 = "raw_data/{sample}_R2.fastq.gz"
    output:
        r1 = "results/trimmed/{sample}_R1_val_1.fq.gz",
        r2 = "results/trimmed/{sample}_R2_val_2.fq.gz"
    threads: 4
    conda: "envs/trim_galore.yaml"
    shell:
        """
        trim_galore --paired --nextera --cores {threads} \
            --output_dir results/trimmed \
            {input.r1} {input.r2}
        """

rule bowtie2_align:
    input:
        r1 = "results/trimmed/{sample}_R1_val_1.fq.gz",
        r2 = "results/trimmed/{sample}_R2_val_2.fq.gz"
    output:
        bam = temp("results/aligned/{sample}_raw.bam")
    threads: 12
    resources:
        mem_mb = 32000
    conda: "envs/bowtie2.yaml"
    shell:
        """
        bowtie2 --very-sensitive -X 2000 --no-mixed \
            --no-discordant --threads {threads} \
            -x {GENOME} \
            -1 {input.r1} -2 {input.r2} \
        | samtools sort -@ {threads} -O bam \
            -o {output.bam}
        """

rule filter_bam:
    input:
        bam = "results/aligned/{sample}_raw.bam"
    output:
        bam = "results/filtered/{sample}_filtered.bam",
        bai = "results/filtered/{sample}_filtered.bam.bai",
        stats = "results/qc/{sample}_filter_stats.txt"
    params:
        blacklist = BLACKLIST
    threads: 4
    conda: "envs/samtools.yaml"
    shell:
        """
        # Remove mito, low-quality, duplicates, blacklist
        samtools view -h -f 2 -q 30 {input.bam} \
            | grep -v 'chrM' \
            | samtools sort -@ {threads} -O bam \
            | samtools markdup -r -@ {threads} - - \
            | bedtools intersect -v -abam stdin \
                -b {params.blacklist} \
            > {output.bam}
        samtools index {output.bam}
        samtools flagstat {output.bam} > {output.stats}
        """

rule tn5_shift:
    input:
        bam = "results/filtered/{sample}_filtered.bam"
    output:
        bam = "results/shifted/{sample}_shifted.bam",
        bai = "results/shifted/{sample}_shifted.bam.bai"
    conda: "envs/deeptools.yaml"
    shell:
        """
        alignmentSieve -b {input.bam} \
            --ATACshift -o {output.bam}
        samtools index {output.bam}
        """

rule macs2_call_peaks:
    input:
        bam = "results/shifted/{sample}_shifted.bam"
    output:
        peaks = "results/peaks/{sample}_peaks.narrowPeak",
        summits = "results/peaks/{sample}_summits.bed"
    params:
        name = "{sample}",
        genome_size = EFFECTIVE_GENOME
    conda: "envs/macs2.yaml"
    shell:
        """
        macs2 callpeak -t {input.bam} \
            -f BAMPE -g {params.genome_size} \
            -n {params.name} \
            --outdir results/peaks \
            --nomodel --keep-dup all -q 0.01
        """

rule make_bigwig:
    input:
        bam = "results/shifted/{sample}_shifted.bam",
        bai = "results/shifted/{sample}_shifted.bam.bai"
    output:
        bw = "results/bigwig/{sample}.bw"
    threads: 4
    params:
        genome_size = EFFECTIVE_GENOME
    conda: "envs/deeptools.yaml"
    shell:
        """
        bamCoverage -b {input.bam} -o {output.bw} \
            --binSize 10 --normalizeUsing RPGC \
            --effectiveGenomeSize {params.genome_size} \
            -p {threads} --extendReads
        """

rule fragment_size_plot:
    input:
        bam = "results/filtered/{sample}_filtered.bam",
        bai = "results/filtered/{sample}_filtered.bam.bai"
    output:
        pdf = "results/qc/{sample}_fragment_sizes.pdf"
    conda: "envs/deeptools.yaml"
    shell:
        """
        bamPEFragmentSize -b {input.bam} \
            --histogram {output.pdf} \
            --maxFragmentLength 1000
        """

rule homer_motifs:
    input:
        peaks = "results/peaks/{sample}_peaks.narrowPeak"
    output:
        dir = directory("results/motifs/homer_results"),
        html = "results/motifs/homer_results/knownResults.html"
    params:
        genome = GENOME
    threads: 4
    conda: "envs/homer.yaml"
    shell:
        """
        findMotifsGenome.pl {input.peaks} \
            {params.genome} {output.dir} \
            -size 200 -p {threads}
        """

rule diffbind_analysis:
    input:
        peaks = expand(
            "results/peaks/{sample}_peaks.narrowPeak",
            sample=SAMPLES),
        bams = expand(
            "results/shifted/{sample}_shifted.bam",
            sample=SAMPLES)
    output:
        results = "results/diffbind/differential_peaks.tsv"
    conda: "envs/diffbind.yaml"
    script:
        "scripts/run_diffbind.R"

rule multiqc:
    input:
        expand("results/qc/{sample}_filter_stats.txt",
               sample=SAMPLES),
        expand("results/qc/{sample}_fragment_sizes.pdf",
               sample=SAMPLES)
    output:
        "results/multiqc/multiqc_report.html"
    conda: "envs/multiqc.yaml"
    shell:
        """
        multiqc results/ -o results/multiqc --force
        """
\end{lstlisting}

\subsection{Nextflow Workflow}
\label{subsec:chromatin:nextflow}

\begin{lstlisting}[language=Java, caption={Nextflow DSL2 workflow for ATAC-seq analysis.}, label={lst:chromatin:nextflow}]
#!/usr/bin/env nextflow
nextflow.enable.dsl=2

params.reads          = "raw_data/*_R{1,2}.fastq.gz"
params.genome_index   = "reference/genome"
params.genome_fasta   = "reference/genome.fa"
params.blacklist      = "reference/blacklist.bed"
params.effective_size = "2.7e9"
params.outdir         = "results"

process TRIM_ADAPTERS {
    tag "$sample_id"
    cpus 4
    publishDir "${params.outdir}/trimmed", mode: 'copy'

    input:
    tuple val(sample_id), path(reads)

    output:
    tuple val(sample_id), path("*_val_{1,2}.fq.gz"),
                                              emit: trimmed
    path "*_trimming_report.txt",             emit: reports

    script:
    """
    trim_galore --paired --nextera \
        --cores ${task.cpus} \
        ${reads[0]} ${reads[1]}
    """
}

process BOWTIE2_ALIGN {
    tag "$sample_id"
    cpus 12
    memory '32 GB'

    input:
    tuple val(sample_id), path(reads)

    output:
    tuple val(sample_id), path("${sample_id}_raw.bam"),
                                              emit: bam

    script:
    """
    bowtie2 --very-sensitive -X 2000 \
        --no-mixed --no-discordant \
        --threads ${task.cpus} \
        -x ${params.genome_index} \
        -1 ${reads[0]} -2 ${reads[1]} \
    | samtools sort -@ ${task.cpus} -O bam \
        -o ${sample_id}_raw.bam
    """
}

process FILTER_BAM {
    tag "$sample_id"
    cpus 4
    publishDir "${params.outdir}/filtered", mode: 'copy'

    input:
    tuple val(sample_id), path(bam)

    output:
    tuple val(sample_id),
          path("${sample_id}_filtered.bam"),
          path("${sample_id}_filtered.bam.bai"),
                                              emit: bam
    path "${sample_id}_flagstat.txt",         emit: stats

    script:
    """
    samtools view -h -f 2 -q 30 ${bam} \
        | grep -v 'chrM' \
        | samtools sort -@ ${task.cpus} -O bam \
        | samtools markdup -r -@ ${task.cpus} - - \
        | bedtools intersect -v -abam stdin \
            -b ${params.blacklist} \
        > ${sample_id}_filtered.bam

    samtools index ${sample_id}_filtered.bam
    samtools flagstat ${sample_id}_filtered.bam \
        > ${sample_id}_flagstat.txt
    """
}

process TN5_SHIFT {
    tag "$sample_id"
    publishDir "${params.outdir}/shifted", mode: 'copy'

    input:
    tuple val(sample_id), path(bam), path(bai)

    output:
    tuple val(sample_id),
          path("${sample_id}_shifted.bam"),
          path("${sample_id}_shifted.bam.bai"),
                                              emit: bam

    script:
    """
    alignmentSieve -b ${bam} --ATACshift \
        -o ${sample_id}_shifted.bam
    samtools index ${sample_id}_shifted.bam
    """
}

process CALL_PEAKS {
    tag "$sample_id"
    publishDir "${params.outdir}/peaks", mode: 'copy'

    input:
    tuple val(sample_id), path(bam), path(bai)

    output:
    tuple val(sample_id),
          path("${sample_id}_peaks.narrowPeak"),
                                              emit: peaks

    script:
    """
    macs2 callpeak -t ${bam} \
        -f BAMPE -g ${params.effective_size} \
        -n ${sample_id} \
        --nomodel --keep-dup all -q 0.01
    """
}

process BIGWIG {
    tag "$sample_id"
    cpus 4
    publishDir "${params.outdir}/bigwig", mode: 'copy'

    input:
    tuple val(sample_id), path(bam), path(bai)

    output:
    path "${sample_id}.bw"

    script:
    """
    bamCoverage -b ${bam} -o ${sample_id}.bw \
        --binSize 10 --normalizeUsing RPGC \
        --effectiveGenomeSize ${params.effective_size} \
        -p ${task.cpus} --extendReads
    """
}

process FOOTPRINTING {
    tag "$sample_id"
    cpus 4
    publishDir "${params.outdir}/footprinting",
               mode: 'copy'

    input:
    tuple val(sample_id), path(bam), path(bai)
    tuple val(sample_id), path(peaks)

    output:
    path "TOBIAS_output/*"

    script:
    """
    TOBIAS ATACorrect -b ${bam} \
        -g ${params.genome_fasta} \
        -p ${peaks} --cores ${task.cpus} \
        --outdir TOBIAS_output

    TOBIAS FootprintScores \
        -s TOBIAS_output/*_corrected.bw \
        -r ${peaks} --cores ${task.cpus} \
        --outdir TOBIAS_output

    TOBIAS BINDetect \
        -m motif_database.jaspar \
        -s TOBIAS_output/*_footprints.bw \
        -b ${peaks} \
        --genome ${params.genome_fasta} \
        --cores ${task.cpus} \
        --outdir TOBIAS_output/bindetect
    """
}

// Main workflow
workflow {
    reads_ch = Channel
        .fromFilePairs(params.reads,
                       checkIfExists: true)

    TRIM_ADAPTERS(reads_ch)
    BOWTIE2_ALIGN(TRIM_ADAPTERS.out.trimmed)
    FILTER_BAM(BOWTIE2_ALIGN.out.bam)
    TN5_SHIFT(FILTER_BAM.out.bam)
    CALL_PEAKS(TN5_SHIFT.out.bam)
    BIGWIG(TN5_SHIFT.out.bam)
    FOOTPRINTING(TN5_SHIFT.out.bam, CALL_PEAKS.out.peaks)
}
\end{lstlisting}

\begin{tipbox}
The nf-core community maintains \software{nf-core/atacseq}, a production-grade Nextflow pipeline that handles all steps from raw reads through peak calling, differential analysis, and motif enrichment. It supports multiple genome builds, includes comprehensive QC metrics (TSS enrichment, FRiP, fragment distributions), and works on institutional HPC clusters and cloud platforms out of the box.
\end{tipbox}

\subsection{Differential Accessibility Analysis}
\label{subsec:chromatin:differential}

To identify regions with significantly different accessibility between conditions:

\begin{enumerate}[itemsep=2pt]
  \item \textbf{Create a consensus peak set:} Merge peaks from all samples (e.g., using \texttt{bedtools merge}), then count reads in each consensus peak for each sample using \software{featureCounts} or \software{bedtools multicov}.
  \item \textbf{Normalise and test:} Use DESeq2 or edgeR on the count matrix to identify differentially accessible regions, or use the \software{DiffBind} R package, which wraps DESeq2/edgeR and provides ATAC-seq-specific functionality including affinity-based analysis.
  \item \textbf{Annotate:} Use \software{ChIPseeker} or HOMER \texttt{annotatePeaks.pl} to assign peaks to nearest genes and genomic features (promoter, intron, intergenic, etc.).
\end{enumerate}

\subsection{Motif Analysis and Footprinting}
\label{subsec:chromatin:motifs}

\textbf{Motif enrichment analysis} identifies transcription factor binding motifs over-represented in accessible regions. Key tools:

\begin{itemize}[itemsep=2pt]
  \item \textbf{HOMER:} \texttt{findMotifsGenome.pl} performs both known-motif enrichment and de novo motif discovery against a background set.
  \item \textbf{chromVAR:} Calculates per-cell TF motif accessibility deviation scores. Ideal for scATAC-seq and for comparing motif accessibility across conditions in bulk data.
  \item \textbf{MEME Suite:} Includes MEME (de novo motif finding), FIMO (individual motif scanning), and AME (motif enrichment testing).
\end{itemize}

\textbf{TF footprinting} goes beyond motif enrichment to detect evidence of actual TF occupancy. When a TF is bound to DNA, it protects a small region from Tn5 cutting, creating a ``footprint'' (a local dip in the otherwise high accessibility signal):

\begin{itemize}[itemsep=2pt]
  \item \textbf{TOBIAS:} State-of-the-art ATAC-seq footprinting tool. Corrects Tn5 sequence bias, calculates footprint scores at motif instances, and uses BINDetect to identify differential TF binding between conditions.
  \item \textbf{HINT-ATAC:} Footprinting framework specifically designed for ATAC-seq data, with built-in bias correction models.
\end{itemize}

\begin{warningbox}
Footprinting analysis requires substantial sequencing depth ($>$100 million uniquely mapped reads per sample) and careful Tn5 bias correction. At typical ATAC-seq depths (50--75 million reads), footprinting is often underpowered for all but the most abundantly expressed TFs. Motif enrichment analysis is more robust at standard sequencing depths and should be the primary analysis for most experiments.
\end{warningbox}

\subsection{TF Footprinting Scores: The TOBIAS Framework}
\label{subsec:chromatin:tobias-framework}

\software{TOBIAS} (Transcription factor Occupancy prediction By Investigation of ATAC-seq Signal) provides a comprehensive framework for detecting transcription factor footprints from ATAC-seq data. The method proceeds through three stages: Tn5 bias correction, footprint scoring, and differential binding detection.

\subsubsection{Tn5 Bias Correction (ATACorrect)}

The Tn5 transposase has an intrinsic sequence preference for insertion, which creates a non-uniform background of cut sites even in the absence of TF binding. If uncorrected, this bias can mimic or mask genuine footprints. TOBIAS corrects for Tn5 bias by:

\begin{enumerate}[itemsep=2pt]
  \item Learning the \textbf{k-mer insertion preference} of Tn5 from the observed data. For each k-mer context (typically 6-mers) surrounding an insertion site, TOBIAS estimates the expected insertion frequency relative to the genomic average.
  \item Computing an \textbf{expected insertion profile} for each genomic position based on the local sequence context and the learned Tn5 bias model.
  \item Subtracting the expected (bias) profile from the observed insertion profile to produce a \textbf{bias-corrected insertion signal}. This corrected signal reflects chromatin accessibility and TF occupancy rather than enzyme preference.
\end{enumerate}

\subsubsection{Footprint Scoring (FootprintScores)}

After bias correction, TOBIAS computes a \textbf{footprint score} at each potential TF binding site (defined by motif occurrences from databases such as JASPAR):

\begin{itemize}[itemsep=2pt]
  \item At a motif site, TOBIAS examines the local pattern of Tn5 insertions. A bound TF creates a characteristic profile: \textbf{high insertion flanking the motif} (accessible DNA on either side of the TF) and \textbf{low insertion within the motif} (the TF physically blocks Tn5).
  \item The protection score quantifies this pattern:
\end{itemize}

\begin{equation}
\label{eq:chromatin:footprint-score}
\text{Footprint score} = \overline{S}_{\text{flanking}} - \overline{S}_{\text{motif}}
\end{equation}

where $\overline{S}_{\text{flanking}}$ is the mean bias-corrected insertion signal in the flanking windows (typically $\pm$50--100\basepair{} around the motif) and $\overline{S}_{\text{motif}}$ is the mean signal within the motif footprint. A high footprint score indicates strong TF protection (likely bound); a low or negative score suggests the motif is unoccupied.

\subsubsection{Differential Binding Detection (BINDetect)}

The final stage of TOBIAS compares footprint scores between conditions to identify TFs with \textbf{differential binding}:
\begin{itemize}[itemsep=2pt]
  \item For each TF motif, TOBIAS collects the footprint scores at all motif instances genome-wide.
  \item It fits a mixture model to classify each motif instance as \textbf{bound} or \textbf{unbound} based on the footprint score distribution.
  \item Differential binding between conditions is assessed by comparing the proportion of bound motif instances and the distribution of footprint scores across conditions.
  \item Output includes a volcano-like plot showing each TF's change in binding activity versus significance.
\end{itemize}

\subsection{Nucleosome Positioning from ATAC-seq Fragment Sizes}
\label{subsec:chromatin:nucleosome-positioning}

A unique feature of ATAC-seq is that the \textbf{fragment size distribution} encodes information about nucleosome positioning. Because Tn5 inserts into accessible linker DNA but is blocked by nucleosomes, the distance between two Tn5 insertion events on opposite sides of a nucleosome reflects the nucleosome footprint.

\subsubsection{Fragment Size Classes}

ATAC-seq fragments can be classified into three biologically meaningful categories based on the 147\basepair{} nucleosome core particle size:

\begin{table}[htbp]
\centering
\caption[ATAC-seq fragment size classes and their biological interpretation.]{Classification of ATAC-seq fragments by size and their chromatin structural interpretation.}
\label{tab:chromatin:fragment-classes}
\small
\begin{tabularx}{\textwidth}{l c X}
\toprule
\textbf{Fragment Class} & \textbf{Size Range} & \textbf{Biological Interpretation} \\
\midrule
Sub-nucleosomal (NFR) & $<$ 147\basepair{} & Both Tn5 insertions in the same nucleosome-free region; represents open chromatin signal \\
\addlinespace
Mono-nucleosomal & 147--294\basepair{} & Tn5 insertions flanking a single nucleosome; the fragment spans one nucleosome core particle \\
\addlinespace
Di-nucleosomal & 294--441\basepair{} & Fragment spans two nucleosomes plus the intervening linker DNA \\
\addlinespace
Multi-nucleosomal & $>$ 441\basepair{} & Fragment spans three or more nucleosomes; rare in well-prepared libraries \\
\bottomrule
\end{tabularx}
\end{table}

\subsubsection{The 147\basepair{} Periodicity}

The characteristic ``nucleosomal ladder'' in the ATAC-seq fragment size distribution arises from the 147\basepair{} periodicity imposed by the nucleosome core particle:
\begin{itemize}[itemsep=2pt]
  \item \textbf{Peak at $\sim$200\basepair{}:} Mono-nucleosomal fragments (147\basepair{} of wrapped DNA + $\sim$50\basepair{} of flanking linker DNA on each side).
  \item \textbf{Peak at $\sim$400\basepair{}:} Di-nucleosomal fragments (two nucleosomes + linker).
  \item \textbf{Peak at $\sim$600\basepair{}:} Tri-nucleosomal fragments (three nucleosomes + linkers).
  \item The exact peak positions depend on the linker DNA length, which varies between species, cell types, and genomic regions (typically 20--80\basepair{} in mammals).
\end{itemize}

\subsubsection{Inferring Nucleosome Occupancy}

By separating fragments into sub-nucleosomal and nucleosomal size classes, one can generate distinct genome-wide signal tracks:

\begin{enumerate}[itemsep=2pt]
  \item \textbf{NFR track} (fragments $< 150$\basepair{}): represents the accessibility signal. This is the standard ATAC-seq open-chromatin signal used for peak calling.
  \item \textbf{Nucleosome track} (fragments 150--300\basepair{}): the midpoints of mono-nucleosomal fragments mark nucleosome dyad positions. A coverage track of these midpoints reveals nucleosome positioning patterns, including:
  \begin{itemize}[itemsep=1pt]
    \item Well-positioned $+1$ and $-1$ nucleosomes flanking active TSSs.
    \item Nucleosome-depleted regions (NDRs) at active regulatory elements.
    \item Phased nucleosome arrays downstream of TSSs.
  \end{itemize}
  \item \textbf{Combined analysis:} Plotting the NFR signal and nucleosome occupancy together around genomic features (e.g., TSSs, TF motifs) provides a comprehensive view of the local chromatin architecture.
\end{enumerate}

\begin{keyconcept}[title={ATAC-seq as a Dual Assay}]
ATAC-seq simultaneously provides two complementary readouts: (1) \textbf{chromatin accessibility} from short, sub-nucleosomal fragments ($< 147$\basepair{}), and (2) \textbf{nucleosome positioning} from mono-nucleosomal fragments ($\sim$147--294\basepair{}). This duality is unique among accessibility assays and enables joint analysis of regulatory element activity and nucleosome architecture from a single experiment. Tools such as \software{NucleoATAC} and \software{HMMRATAC} exploit this property to jointly call accessible regions and position nucleosomes from ATAC-seq data.
\end{keyconcept}

%% ============================================================
\section{Experimental Design Considerations}
\label{sec:chromatin:design}

\begin{keyconcept}[title={ATAC-seq Experimental Design Checklist}]
\begin{itemize}[itemsep=2pt]
  \item \textbf{Cell number:} 50,000--100,000 cells per reaction (Omni-ATAC). Fresh cells preferred; flash-frozen nuclei are acceptable.
  \item \textbf{Replicates:} At least 2 biological replicates for peak calling, 3+ for differential analysis.
  \item \textbf{Sequencing:} Paired-end 50 or 75\basepair{}; 50--75M read pairs per sample for bulk ATAC-seq.
  \item \textbf{Controls:} No input control is needed (unlike ChIP-seq), but include appropriate biological controls (e.g., untreated vs.\ treated).
  \item \textbf{Avoid:} over-tagmentation (destroys signal by fragmenting all chromatin), under-tagmentation (low library complexity), and over-amplification (excessive PCR duplicates).
  \item \textbf{Fresh Tn5:} Enzyme activity degrades over time; use aliquots and avoid repeated freeze-thaw cycles.
\end{itemize}
\end{keyconcept}

%% ============================================================
\section{Chapter Summary}
\label{sec:chromatin:summary}

Chromatin accessibility profiling reveals the regulatory landscape of cells by mapping open chromatin regions genome-wide. ATAC-seq has become the dominant method due to its simplicity (50,000 cells, approximately one hour of bench work), and the Omni-ATAC variant further improves signal quality while reducing mitochondrial contamination. Single-cell ATAC-seq extends this approach to individual cells, albeit with extreme data sparsity that requires specialised computational approaches such as latent semantic indexing, topic modelling, and gene activity score estimation. The bioinformatics pipeline for ATAC-seq---from alignment and filtering through peak calling, differential accessibility, motif enrichment, and footprinting---provides a comprehensive view of the transcription factor binding landscape. Integration of accessibility data with gene expression (\acrshort{scRNA-seq}), histone modifications (\acrshort{ChIP-seq}/CUT\&Tag), DNA methylation, and 3D genome architecture (Hi-C) data produces a multidimensional view of gene regulation that is far richer than any single assay alone.
